\documentclass[10pt]{article}
\usepackage[fr]{acta2}   % [fr] or [en]
\renewcommand{\actaedno}{319}\renewcommand{\actacode}{AP-2026-319-MA-AGR-2-S}
\newcommand{\jname}{Prénom Nom}\newcommand{\jmail}{prenom.nom@email.ma}
\newcommand{\jlink}{linkedin.com/in/votre-profil}\newcommand{\jtel}{+212 6 00 00 00 00}
\begin{document}
\atitle{\sSolution\ // \sSector: Agroalimentaire // \sCity: Fès-Meknès}{Dimensionnement d'une ligne de pompage}
\astatus{\sRealistic}{\sHonest}
\asec{01}{\sProblem}
Pomper $Q=\SI{40}{\meter\cubed\per\hour}$ d'eau sur $\Delta z=\SI{18}{\meter}$ par une conduite DN\,100 de $L=\SI{250}{\meter}$. Déterminer $H$ et la puissance à l'arbre.
\asec{02}{\sAssump}
Écoulement permanent, eau à \SI{20}{\celsius}, acier commercial ($\varepsilon=\SI{0.045}{\milli\meter}$), $\sum K=6{,}5$, rendement $\eta=0{,}68$.
\asec{03}{\sCalc}
\begin{minipage}[t]{0.5\textwidth}\vspace{0pt}\small
\begin{align*}
v&=\tfrac{4Q}{\pi D^{2}}=\SI{1.28}{\meter\per\second}\\
\mathrm{Re}&=\tfrac{vD}{\nu}=1{,}34\!\times\!10^{5}\Rightarrow\text{turbulent}\\
f&=0{,}0194\ \ (\text{Swamee--Jain})\\
h_f&=f\tfrac{L}{D}\tfrac{v^{2}}{2g}=\SI{3.87}{\meter},\ h_m=\SI{0.55}{\meter}\\
H&=18+3{,}87+0{,}55=\SI{22.4}{\meter}\\
P&=\tfrac{\rho g Q H}{\eta}=\SI{3.59}{\kilo\watt}
\end{align*}\end{minipage}\hfill
\begin{minipage}[t]{0.46\textwidth}\vspace{0pt}\centering
\begin{tikzpicture}
\begin{axis}[width=\linewidth,height=5.2cm,xlabel={$Q$ (m$^3$/h)},ylabel={$H$ (m)},axis lines=left,
 axis line style={ink},tick style={ink},grid=major,grid style={ink!25,dotted},tick label style={font=\tiny},label style={font=\scriptsize},ymin=16,ymax=29,xmin=0,xmax=60]
 \addplot[very thick,orange,smooth,mark=*,mark size=1.2pt]coordinates{(0,18)(10,18.33)(20,19.2)(30,20.56)(40,22.41)(50,24.75)(60,27.56)};
 \draw[ink,dashed](axis cs:40,16)--(axis cs:40,22.41)--(axis cs:0,22.41);
 \node[font=\tiny,anchor=south east]at(axis cs:40,22.9){(40 ; 22,4)};
\end{axis}\end{tikzpicture}\\[-1pt]{\monofont\tiny\MakeUppercase{\sFigure}\ :\ H(Q)}
\end{minipage}
\asec{04}{\sResult}
\colorbox{ink}{\parbox{\dimexpr\textwidth-2\fboxsep}{\color{paper}\small\textbf{$H\approx\SI{22.4}{\meter}$,\ $P_{\text{arbre}}\approx\SI{3.6}{\kilo\watt}$.} Une pompe de \SIrange{4}{5.5}{\kilo\watt} avec variateur de vitesse couvre le point de fonctionnement avec marge ; le frottement ne pèse que 17\,\% de $H$, le gain vient du choix du point de fonctionnement.}}
\vfill
{\monofont\color{orange}\scriptsize\MakeUppercase{\sWho}}\ \ {\bfseries\jname}\\
{\monofont\scriptsize\MakeUppercase{\sContact}\,:\ \jmail\ //\ \jtel\ //\ \jlink}
\end{document}
